\section{Control reset-free y gestión de la memoria}
\subsection{Debilitamiento del supuesto de reinicio}
El RL tradicional depende de reinicios hechos por una persona o por el simulator. Chelsea propone «que la propia policy haga el reset»: al terminar cada episodio de la tarea se activa otra policy que devuelve el entorno a un nuevo conjunto inicial, de modo que el robot pueda operar de forma continua.

\begin{figure}[H]
\centering
\includegraphics[width=0.88\textwidth]{slides-images/slide-04.jpg}
\caption{La slide 4 explica cómo los controladores hacia adelante y hacia atrás alternan entre ejecutar la tarea y hacer el reset.}
\end{figure}

\noindent La timeline de la figura representa el proceso de recuperación después de una ejecución de la tarea: cuando la policy de la tarea (al centro) alcanza el goal, entrega el current state a la reset policy. Al ejecutarse, la reset policy lleva el system de vuelta a la base manifold y sigue proporcionando puntos de partida para la tarea.

\begin{importantbox}{Mecanismo del Forward-Backward Controller}
\begin{itemize}
\item $\pi_f$: explora, a partir del initial state, una task-specific trajectory;\\
\item $\pi_b$: recibe el state actual e intenta volver a la initial manifold;\\
\item Se entrenan ambas por turnos; la tasa de éxito de $\pi_b$ define el punto de partida del nuevo episode.
\end{itemize}
\end{importantbox}

\subsection{Episodic buffer persistente}
El enfoque reset-free requiere un long-lived buffer que recuerde los «failed states más recientes», para que la política siga aprendiendo a hacer el reinicio. Chelsea recomienda mantener un prioritized replay buffer e incorporar además \textbf{time-aware sampling}, que da prioridad a reproducir los datos generados después del reset más reciente.

\begin{knowledgebox}{Time-Aware Replay}
\begin{itemize}
\item Se añade un timestamp a cada transition;\\
\item Durante el entrenamiento se usa $w_t \propto \exp(-\lambda (t_{now}-t))$ para aproximar la latest priority;\\
\item Así se garantiza que los datos nuevos se aprovechen plenamente, sin dejar de conservar algunos rare events.
\end{itemize}
\end{knowledgebox}

\subsection{Deshacer experiencia y diversidad del entorno}
El ponente señala además que un sistema autónomo debe tener la capacidad de \textbf{deshacer experiencia}, es decir, que el modelo retroceda sobre las trayectorias en las que hizo fail y pruebe nuevas combinaciones de estrategias. Esto implica planificar de forma dinámica varias reset-policies para cubrir distintas perturbaciones físicas.
\begin{warningbox}{Fail-fast but recover-fast}
Se alienta a la política a probar con rapidez varias combinaciones de acciones y, si el control actual falla, a activar de inmediato la reset policy; al mismo tiempo, se registran los modos de falla para la future policy adaptation.
\end{warningbox}

\subsection{Resumen de la sección}
El marco reset-free usa forward-backward controllers, un time-aware buffer y fast recovery para lograr una operación continua sin depender de resets manuales.
\newpage
